% =====================================================================
%  ĐỀ KIỂM TRA ĐẦU VÀO — CSCA TOÁN
%  Rút gọn từ đề thi chính thức CSCA 20260315 (48 câu → 20 câu)
% =====================================================================
\documentclass[11pt, a4paper]{article}
% =====================================================================
%  CAE SHANGHAI — GÓI LỆNH DÙNG CHUNG CHO TÀI LIỆU LUYỆN THI CSCA TOÁN
%  Biên dịch bằng XeLaTeX. Mọi file .tex trong buổi học \input file này.
% =====================================================================

% ---------- NGÔN NGỮ & PHÔNG CHỮ ----------
\usepackage{fontspec}
\usepackage{polyglossia}
\setmainlanguage{vietnamese}
\setotherlanguage{english}

% Chữ Latin (Anh + Việt): TeX Gyre Termes — bản clone Times New Roman,
% đi kèm TeX Live, hỗ trợ đầy đủ dấu tiếng Việt.
\setmainfont{TeX Gyre Termes}
\setsansfont{TeX Gyre Heros}

% Chữ Hán: ưu tiên Source Han Serif SC (思源宋体) — chuẩn học thuật Trung Quốc.
% Nếu máy chưa cài, tự động dùng Songti SC (宋体, có sẵn trên macOS).
% Cả hai đều thuộc họ 宋体 nên diện mạo tài liệu không đổi.
\IfFontExistsTF{Source Han Serif SC}{
  \newcommand{\cjkfamilyname}{Source Han Serif SC}
}{
  \IfFontExistsTF{Songti SC}{
    \newcommand{\cjkfamilyname}{Songti SC}
  }{
    \newcommand{\cjkfamilyname}{Noto Serif CJK SC}
  }
}

% xeCJK tự chuyển font khi gặp ký tự Hán, không cần bọc \cjk{} thủ công.
% Đây là điều kiện bắt buộc: font Latin (TeX Gyre Termes) không có glyph chữ Hán,
% nếu không khai báo xeCJK thì mọi chữ Hán viết trần sẽ bị nuốt khỏi bản in.
\usepackage{xeCJK}
\setCJKfamilyfont{caehan}{\cjkfamilyname}[Scale=0.95]
\setCJKmainfont{\cjkfamilyname}[Scale=0.95]
\setCJKsansfont{\cjkfamilyname}[Scale=0.95]
\setCJKmonofont{\cjkfamilyname}[Scale=0.95]

% \cjk{} ép dùng font Hán trong ngữ cảnh cần thiết (ví dụ trong tiêu đề in đậm
% hoặc trong bảng có định dạng riêng). Dùng \CJKfamily chứ không phải tên font,
% vì tên font chỉ là văn bản — sẽ in ra chữ thay vì đổi font.
\newcommand{\cjk}[1]{{\CJKfamily{caehan}#1}}

% ---------- TOÁN HỌC ----------
\usepackage{amsmath, amssymb, amsthm, mathtools}
\usepackage{unicode-math}
\setmathfont{TeX Gyre Termes Math}

% ---------- BỐ CỤC TRANG ----------
% Lề phải chừa vùng an toàn của letterhead:
%   banner trên chiếm 0–19mm  → top tối thiểu 2,2cm
%   banner dưới chiếm 281–297mm → bottom tối thiểu 2,2cm
% Lề trên 2,6cm (26mm) để tiêu đề đầu trang không đè banner (banner hết ở 19mm).
\usepackage[a4paper, top=2.6cm, bottom=2.4cm, left=2.5cm, right=2.0cm]{geometry}
\usepackage{fancyhdr}
\usepackage{titlesec}
\usepackage{enumitem}
\usepackage{graphicx}
\usepackage{booktabs, array, longtable}
\usepackage{xcolor}
\usepackage[most]{tcolorbox}
\usepackage[colorlinks=true, linkcolor=cae_blue, citecolor=cae_blue, urlcolor=cae_blue]{hyperref}

% ---------- BẢNG MÀU THƯƠNG HIỆU ----------
\definecolor{cae_blue}{RGB}{0, 51, 102}
\definecolor{cae_gold}{RGB}{196, 149, 60}
\definecolor{box_bg}{RGB}{245, 247, 250}
\definecolor{box_line}{RGB}{0, 51, 102}

% ---------- TIÊU ĐỀ MỤC ----------
% Đệm đầu trang: banner letterhead chiếm 0–19mm, mà lề trên là 22mm.
% \section* đầu trang (sau \clearpage) sát lề nên phải chừa thêm, nếu không
% sẽ đè lên banner.
\titleformat{\section}
  {\normalfont\large\bfseries\sffamily\color{cae_blue}}{\thesection.}{0.6em}{}
\titleformat{\subsection}
  {\normalfont\normalsize\bfseries\sffamily\color{cae_blue!85!black}}{\thesubsection}{0.6em}{}
\titlespacing*{\section}{0pt}{14pt}{6pt}
\titlespacing*{\subsection}{0pt}{10pt}{4pt}

% Tiêu đề không đánh số ở đầu trang: \section* sau \clearpage nằm sát lề trên
% nên đè lên banner letterhead. Dùng \vspace* ngay TRƯỚC \section* không đủ,
% vì titlesec chèn khoảng cách riêng. Cách chắc chắn là đệm bằng \null.
\newcommand{\sectiondautrang}[1]{\section*{#1}}

% Dạng bài dùng \subsection*{} nên cần định dạng riêng, đậm và có gạch chân màu
\titleformat{\subsubsection}
  {\normalfont\normalsize\bfseries\sffamily\color{cae_blue!85!black}}{\thesubsubsection}{0.6em}{}

% ---------- LETTERHEAD (ảnh nền toàn trang) ----------
% Ảnh letterhead của CAE Shanghai: banner xanh trên/dưới kèm logo chìm giữa.
% Tỉ lệ ảnh 1414x2000 = 0.7070, khớp A4 (0.7071) nên chèn thẳng không méo.
% Đặt ở lớp nền dưới cùng để chữ không bị ảnh che.
% Dùng \AddToShipoutPictureBG (không có dấu *) để áp dụng cho MỌI trang.
\usepackage{eso-pic}
\AddToShipoutPictureBG{%
  \AtPageLowerLeft{%
    \includegraphics[width=\paperwidth,height=\paperheight]{images/letterhead.png}%
  }%
}

% ---------- ĐẦU / CHÂN TRANG ----------
% Letterhead đã có banner trên và dưới, nên bỏ dòng chạy để tránh chồng chữ.
% Số trang đặt trong vùng nội dung, sát trên banner dưới, lề phải.
\pagestyle{fancy}
\setlength{\headheight}{0pt}
\fancyhf{}
\renewcommand{\headrulewidth}{0pt}
\renewcommand{\footrulewidth}{0pt}
\fancyfoot[R]{\raisebox{0.9cm}{\small\color{cae_blue}\thepage}}

% ---------- MÔI TRƯỜNG NỘI DUNG ----------
\newtcolorbox{dinhnghia}[1][]{
  colback=box_bg, colframe=box_line, boxrule=0.6pt, arc=1pt,
  fonttitle=\bfseries, title={Định nghĩa #1}, breakable}
\newtcolorbox{dinhly}[1][]{
  colback=box_bg, colframe=cae_gold, boxrule=0.6pt, arc=1pt,
  fonttitle=\bfseries, title={Tính chất #1}, breakable}
\newtcolorbox{vidu}[1][]{
  colback=white, colframe=black!40, boxrule=0.4pt, arc=1pt,
  fonttitle=\bfseries, title={Ví dụ #1}, breakable}
\newtcolorbox{baitap}[1][]{
  colback=white, colframe=cae_blue!50, boxrule=0.4pt, arc=1pt,
  fonttitle=\bfseries, title={Câu #1}, breakable}

% Ô đáp án: viền xanh đậm hơn, dùng trong file đáp án
\newtcolorbox{dapan}[1][]{
  colback=box_bg, colframe=cae_blue, boxrule=0.6pt, arc=1pt,
  fonttitle=\bfseries, title={Câu #1}, breakable}

% ---------- TÊN CHỨNG MINH ----------
\renewcommand{\proofname}{\normalfont\bfseries Chứng minh}

% ---------- TRANG BÌA GỌN ----------
% #1 = loại tài liệu (TÀI LIỆU CHUẨN BỊ, TÀI LIỆU BUỔI HỌC, ...)
% #2 = số buổi
% #3 = tên buổi tiếng Việt
% #4 = tên buổi tiếng Hán
% #5 = module và tỉ trọng
\newcommand{\trangbia}[5]{%
  \begin{center}
    \vspace*{0.6cm}
    {\large\bfseries\color{cae_blue} CAE SHANGHAI}\\[0.15cm]
    {\footnotesize Tài liệu luyện thi CSCA · Môn Toán}\\[0.5cm]
    {\color{cae_gold}\rule{0.55\linewidth}{0.8pt}}\\[0.35cm]
    {\small\bfseries\color{cae_blue!80!black} #1}\\[0.35cm]
    {\Large\bfseries\color{cae_blue} BUỔI #2 — #3}\\[0.2cm]
    {\large\color{cae_blue!85!black} \cjk{#4}}\\[0.25cm]
    {\small #5}\\[0.35cm]
    {\color{cae_gold}\rule{0.55\linewidth}{0.8pt}}\\[0.3cm]
    {\footnotesize Biên soạn: Hoàng Dương, CAE Shanghai \quad·\quad \today}
  \end{center}
  \vspace{0.3cm}
}

% ---------- TRANG BÌA CHO ĐỀ KIỂM TRA ----------
% Khác \trangbia ở chỗ không có chữ "BUỔI N" — dùng cho đề thi, không phải bài giảng.
% #1 = loại tài liệu, #2 = tên kỳ thi (Việt), #3 = tên kỳ thi (Hán), #4 = thông số
\newcommand{\biade}[4]{%
  \begin{center}
    \vspace*{0.6cm}
    {\large\bfseries\color{cae_blue} CAE SHANGHAI}\\[0.15cm]
    {\footnotesize Tài liệu luyện thi CSCA · Môn Toán}\\[0.5cm]
    {\color{cae_gold}\rule{0.55\linewidth}{0.8pt}}\\[0.35cm]
    {\small\bfseries\color{cae_blue!80!black} #1}\\[0.35cm]
    {\Large\bfseries\color{cae_blue} #2}\\[0.2cm]
    {\large\color{cae_blue!85!black} \cjk{#3}}\\[0.25cm]
    {\small #4}\\[0.35cm]
    {\color{cae_gold}\rule{0.55\linewidth}{0.8pt}}\\[0.3cm]
    {\footnotesize Biên soạn: Hoàng Dương, CAE Shanghai \quad·\quad \today}
  \end{center}
  \vspace{0.3cm}
}

% ---------- TIỆN ÍCH ----------
% Nhãn từ khóa in đậm ở đầu đoạn, dùng cho Từ vựng / Từ khóa nhận dạng / Chú ý
\newcommand{\nhan}[1]{\textbf{#1.}\ }

% Bảng gọn, không đánh số
\newcommand{\bang}[1]{%
  \begin{center}#1\end{center}}

% Đánh dấu ô trống cho học sinh điền
\newcommand{\chotrong}{\makebox[3.2cm]{\dotfill}}


\begin{document}

\biade{ĐỀ KIỂM TRA ĐẦU VÀO}{CSCA Toán}{数学能力测试}{20 câu \quad·\quad 30 phút \quad·\quad không dùng máy tính}

\vspace{0.2cm}
\noindent\emph{Đề gồm 20 câu trắc nghiệm, mỗi câu chọn một đáp án đúng. Thời gian 30 phút, trung bình 90 giây mỗi câu. Không được dùng máy tính.}

\vspace{0.1cm}
% =====================================================================
\section*{Phần 1 — Tập hợp và bất đẳng thức (2 câu)}
\addcontentsline{toc}{section}{Phần 1 — Tập hợp và bất đẳng thức}

\begin{baitap}[1]
设集合 $A=\{x \mid -2<x\le 3\}$。下列哪个说法是正确的？（ ）\\[4pt]
A. $2\notin A$ \quad B. $-3\in A$ \quad C. $3\in A$ \quad D. $-2\in A$
\end{baitap}

\begin{baitap}[2]
不等式 $x^2-x-6>0$ 的解集是（ ）\\[4pt]
A. $\{x \mid -2<x<3\}$ \quad B. $\{x \mid x<-2 \text{ 或 } x>3\}$ \quad C. $\{x \mid x<-2\}$ \quad D. $\{x \mid x>3\}$
\end{baitap}

% =====================================================================
\section*{Phần 2 — Hàm số và dãy số (10 câu)}
\addcontentsline{toc}{section}{Phần 2 — Hàm số và dãy số}

\begin{baitap}[3]
设 $\alpha=\dfrac{5\pi}{6}$，则 $\cos\alpha=$（ ）\\[4pt]
A. $-\dfrac{1}{2}$ \quad B. $\dfrac{\sqrt{3}}{2}$ \quad C. $\dfrac{1}{2}$ \quad D. $-\dfrac{\sqrt{3}}{2}$
\end{baitap}

\begin{baitap}[4]
在其定义域上，函数 $f(x)=\ln(x^2)$ 是（ ）\\[4pt]
A. 奇函数 \quad B. 奇偶性不确定 \quad C. 既不是奇函数也不是偶函数 \quad D. 偶函数
\end{baitap}

\begin{baitap}[5]
函数 $f(x)=\sin x$ 单调递增的区间是（ ）\\[4pt]
A. $[-\pi,\pi]$ \quad B. $[0,\pi]$ \quad C. $\left[-\dfrac{\pi}{2},\dfrac{\pi}{2}\right]$ \quad D. $(-\infty,+\infty)$
\end{baitap}

\begin{baitap}[6]
设 $\{a_n\}$ 是等差数列，首项 $a_1=2$，公差 $d=3$。则 $a_{100}=$（ ）\\[4pt]
A. 298 \quad B. 302 \quad C. 305 \quad D. 299
\end{baitap}

\begin{baitap}[7]
若指数函数 $y=a^x$ 在 $\mathbb{R}$ 上单调递增，则 $a$ 的取值范围是（ ）\\[4pt]
A. $a>1$ \quad B. $a<0$ 或 $a>1$ \quad C. $0<a<1$ \quad D. $a>0$
\end{baitap}

\begin{baitap}[8]
设 $a=\lg 4$，$b=\lg 5$。则 $a+2b=$（ ）\\[4pt]
A. 100 \quad B. 4 \quad C. 2 \quad D. 5
\end{baitap}

\begin{baitap}[9]
函数 $f(x)=3\cos\left(4x-\dfrac{\pi}{3}\right)$ 的最小正周期是（ ）\\[4pt]
A. $\dfrac{\pi}{3}$ \quad B. $\dfrac{\pi}{2}$ \quad C. $\dfrac{\pi}{4}$ \quad D. $2\pi$
\end{baitap}

\begin{baitap}[10]
已知 $\cos\alpha=-\dfrac{1}{2}$，$\alpha\in\left(\dfrac{\pi}{2},\pi\right)$，则 $\sin\dfrac{\alpha}{2}$ 的值为（ ）\\[4pt]
A. $\dfrac{\sqrt{3}}{2}$ \quad B. $-\dfrac{1}{2}$ \quad C. $-\dfrac{\sqrt{3}}{2}$ \quad D. $\dfrac{1}{2}$
\end{baitap}

\begin{baitap}[11]
若 $\cos\alpha=\dfrac{1}{3}$ 且 $-\dfrac{\pi}{2}<\alpha<0$，则 $\tan\dfrac{\alpha}{2}=$（ ）\\[4pt]
A. $\dfrac{\sqrt{2}}{2}$ \quad B. $-\sqrt{2}$ \quad C. $-\dfrac{\sqrt{2}}{2}$ \quad D. $\sqrt{2}$
\end{baitap}

\begin{baitap}[12]
设 $f(x)=\log_5(2x+1)$。则不等式 $f(x)\ge 0$ 的解集是（ ）\\[4pt]
A. $\{x \mid x\ge 1\}$ \quad B. $\{x \mid x>0\}$ \quad C. $\{x \mid x>1\}$ \quad D. $\{x \mid x\ge 0\}$
\end{baitap}

% =====================================================================
\clearpage
\sectiondautrang{Phần 3 — Hình học và đại số (7 câu)}
\addcontentsline{toc}{section}{Phần 3 — Hình học và đại số}

\begin{baitap}[13]
直线 $x+\sqrt{3}y-2=0$ 的倾斜角是（ ）\\[4pt]
A. $60^\circ$ \quad B. $120^\circ$ \quad C. $150^\circ$ \quad D. $30^\circ$
\end{baitap}

\begin{baitap}[14]
已知点 $A(-2,-1)$ 和 $B(a,3)$，且 $|AB|=5$。则 $a$ 的值为（ ）\\[4pt]
A. $1$ 或 $-5$ \quad B. $-1$ 或 $5$ \quad C. $\pm 5$ \quad D. $\pm 1$
\end{baitap}

\begin{baitap}[15]
若直线 $l$ 经过两点 $M(0,-1)$ 和 $N(-2,0)$，则 $l$ 的方程为（ ）\\[4pt]
A. $x+y+2=0$ \quad B. $x+2y+2=0$ \quad C. $2x+y+2=0$ \quad D. $2x+y+1=0$
\end{baitap}

\begin{baitap}[16]
圆心在 $A(-3,2)$，半径为 $3$ 的圆的方程是（ ）\\[4pt]
A. $(x+3)^2+(y-2)^2=3$ \quad B. $(x+3)^2+(y-2)^2=9$ \quad C. $(x-3)^2+(y+2)^2=9$ \quad D. $(x-3)^2+(y+2)^2=3$
\end{baitap}

\begin{baitap}[17]
已知椭圆的长轴长为 $10$，两焦点 $F_1$、$F_2$ 分别在 $(3,0)$ 和 $(-3,0)$。则该椭圆的方程为（ ）\\[4pt]
A. $\dfrac{x^2}{9}+\dfrac{y^2}{25}=1$ \quad B. $\dfrac{x^2}{25}+\dfrac{y^2}{16}=1$ \quad C. $\dfrac{x^2}{16}+\dfrac{y^2}{25}=1$ \quad D. $\dfrac{x^2}{16}+\dfrac{y^2}{9}=1$
\end{baitap}

\begin{baitap}[18]
设向量 $\vec{a}=(-2,3)$，向量 $\vec{b}=(2,1)$。则 $\vec{a}-2\vec{b}=$（ ）\\[4pt]
A. $(-6,-1)$ \quad B. $(6,-1)$ \quad C. $(-6,1)$ \quad D. $(6,1)$
\end{baitap}

\begin{baitap}[19]
若直线 $l$ 垂直于直线 $l_1:x+2y+4=0$，且经过直线 $l_2:x+y+1=0$ 和 $l_3:2x+y-1=0$ 的交点，则 $l$ 的方程为（ ）\\[4pt]
A. $2x-y-7=0$ \quad B. $x+2y-3=0$ \quad C. $x-2y-1=0$ \quad D. $2x+y-5=0$
\end{baitap}

% =====================================================================
\section*{Phần 4 — Xác suất và thống kê (1 câu)}
\addcontentsline{toc}{section}{Phần 4 — Xác suất và thống kê}

\begin{baitap}[20]
一个袋子里有五个大小形状完全相同的小球，包括两个红球（标号 1、2）、两个黑球（标号 1、2）和一个白球（标号 1）。现从中随机抽取两个球（不放回）。求两个球颜色不同的概率，以及至少有一个球标号为 1 的概率，这两个概率分别为（ ）\\[4pt]
A. $\dfrac{4}{5}$，$\dfrac{7}{10}$ \quad B. $\dfrac{3}{5}$，$\dfrac{7}{10}$ \quad C. $\dfrac{3}{5}$，$\dfrac{9}{10}$ \quad D. $\dfrac{4}{5}$，$\dfrac{9}{10}$
\end{baitap}

\vspace{0.5cm}
\begin{center}
\textbf{--- Hết ---}
\end{center}

\vspace{0.3cm}
\noindent\emph{Thời gian làm bài: 30 phút. Nộp bài qua hệ thống học trực tuyến.}

\end{document}
